\documentclass[11pt]{article}
% ===========================================================================
%  Shared preamble for all MPM2D worksheets — input right after \documentclass.
%  (Files beginning with "_" are skipped by mpm2d-worksheet-publish.mjs.)
% ===========================================================================
\usepackage[letterpaper,top=2.2cm,bottom=1.9cm,left=1.6cm,right=1.6cm,headheight=28pt,footskip=24pt]{geometry}
\usepackage{amsmath,amssymb}
\usepackage{cancel}
\usepackage[T1]{fontenc}
\usepackage[utf8]{inputenc}
\usepackage{lmodern}
\usepackage{xcolor}
\usepackage[most]{tcolorbox}
\usepackage{enumitem}
\usepackage{multicol}
\usepackage{fancyhdr}
\usepackage{tikz}
\usetikzlibrary{calc}
\usepackage{pgfplots}
\pgfplotsset{compat=1.18}

\definecolor{exblue}{HTML}{4A90E2}
\definecolor{exbg}{HTML}{E6F3FF}
\definecolor{qorange}{HTML}{E69138}
\definecolor{qbg}{HTML}{FFF7CC}
\definecolor{keygray}{HTML}{475569}
\definecolor{keybg}{HTML}{F1F5F9}
\definecolor{iagreen}{HTML}{14653B}
\definecolor{titleblue}{HTML}{1D4ED8}
% Rotating palette so each worked example gets its own colour.
\definecolor{ex0}{HTML}{2563A0}\definecolor{ex1}{HTML}{3B7D3B}\definecolor{ex2}{HTML}{8A5A00}
\definecolor{ex3}{HTML}{A3327A}\definecolor{ex4}{HTML}{6D28A3}\definecolor{ex5}{HTML}{B08900}
\definecolor{ex6}{HTML}{0E7490}\definecolor{ex7}{HTML}{9A3412}\definecolor{ex8}{HTML}{15803D}

\pagestyle{fancy}
\fancyhf{}
\fancyhead[L]{\footnotesize NAME: \underline{\hspace{3.4cm}}}
\fancyhead[C]{\textbf{Integration Academy}\\[-2pt]{\scriptsize Math Simplified \textperiodcentered{} Success Amplified}}
\fancyhead[R]{\footnotesize MPM2D}
\fancyfoot[L]{\scriptsize dr.merie@IntegrationAcademy.ca}
\fancyfoot[C]{\scriptsize\thepage}
\fancyfoot[R]{\scriptsize IntegrationAcademy.ca}
\renewcommand{\headrulewidth}{1.2pt}
\renewcommand{\footrulewidth}{0.4pt}
\setlength{\parindent}{0pt}
\setlength{\parskip}{4pt}

% Examples are NOT breakable, so each example + its solution stays on one page.
\newtcolorbox{example}[2]{colframe=#1,colback=#1!7,coltitle=white,colbacktitle=#1,fonttitle=\bfseries,title={#2},arc=3pt,boxrule=1pt}
\newtcolorbox{qbox}[1]{colframe=qorange,colback=qbg,coltitle=white,colbacktitle=qorange,fonttitle=\bfseries,title={#1},arc=3pt,breakable,boxrule=1pt}
\newtcolorbox{keybox}{colframe=keygray,colback=keybg,coltitle=white,colbacktitle=keygray,fonttitle=\bfseries,title={$\checkmark$\ Answer Key},arc=3pt,breakable,boxrule=1pt}
\newtcolorbox{recapbox}{colframe=keygray,colback=keybg,coltitle=white,colbacktitle=keygray,fonttitle=\bfseries,title={Question Recap},arc=3pt,breakable,boxrule=1pt}
\newtcolorbox{ideabox}{colframe=titleblue,colback=titleblue!6,coltitle=white,colbacktitle=titleblue,fonttitle=\bfseries,title={The Big Ideas},arc=3pt,breakable,boxrule=1.2pt}
\newtcolorbox{learnbox}{colframe=iagreen,colback=iagreen!5,coltitle=white,colbacktitle=iagreen,fonttitle=\bfseries,title={The Big Picture},arc=3pt,breakable,boxrule=1.2pt}
\newcommand{\lh}[1]{\par\smallskip{\bfseries\color{iagreen}#1}\par\nobreak\smallskip}

\newcommand{\soln}{\par\smallskip\textit{\textbf{Solution.}}\ }
\newcommand{\workspace}[1]{\par\vspace{3pt}\fcolorbox{black!30}{white}{\begin{minipage}[t][#1][t]{\dimexpr\linewidth-2\fboxsep\relax}\ \end{minipage}}\par}
\newcommand{\sech}[1]{\par\vspace{8pt}{\Large\bfseries\color{iagreen}#1}\par\nobreak\vspace{2pt}{\color{black!20}\hrule height1pt}\vspace{6pt}}

% A compact coordinate plot sized for an example box: \eplot{xmin}{xmax}{ymin}{ymax}{body}
\newcommand{\eplot}[5]{\begin{center}\begin{tikzpicture}
\begin{axis}[width=6.8cm,height=4.4cm,axis lines=middle,xlabel={\tiny$x$},ylabel={\tiny$y$},
grid=both,grid style={gray!16},xmin=#1,xmax=#2,ymin=#3,ymax=#4,
xtick distance=2,ytick distance=2,tick label style={font=\tiny},axis line style={-},samples=120]
#5
\end{axis}\end{tikzpicture}\end{center}}

% A sine-curve plot (degrees on x, 0..360): \splot{ymin}{ymax}{body}
\newcommand{\splot}[3]{\begin{center}\begin{tikzpicture}
\begin{axis}[width=8.6cm,height=4.6cm,axis lines=middle,xlabel={\tiny$x\,(^\circ)$},ylabel={\tiny$y$},
grid=both,grid style={gray!16},xmin=0,xmax=360,ymin=#1,ymax=#2,
xtick distance=90,ytick distance=1,tick label style={font=\tiny},axis line style={-},samples=180]
#3
\end{axis}\end{tikzpicture}\end{center}}

% A small coordinate plot: \plot{xmin}{xmax}{ymin}{ymax}{body}
\newcommand{\plot}[5]{\begin{center}\begin{tikzpicture}
\begin{axis}[width=8.4cm,height=6cm,axis lines=middle,xlabel={\scriptsize$x$},ylabel={\scriptsize$y$},
grid=both,grid style={gray!16},xmin=#1,xmax=#2,ymin=#3,ymax=#4,
xtick distance=2,ytick distance=2,tick label style={font=\tiny},axis line style={-}]
#5
\end{axis}\end{tikzpicture}\end{center}}

% Right triangle (angle theta at bottom-left, right angle at bottom-right):
%   \rtri{drawBase}{drawHeight}{oppLabel}{adjLabel}{hypLabel}
\newcommand{\rtri}[5]{\begin{center}\begin{tikzpicture}[scale=0.85]
\coordinate (A) at (0,0);\coordinate (B) at (#1,0);\coordinate (C) at (#1,#2);
\draw[thick,exblue] (A)--(B)--(C)--cycle;
\draw[black] ($(B)+(-0.32,0)$)--($(B)+(-0.32,0.32)$)--($(B)+(0,0.32)$);
\node[below] at ($(A)!0.5!(B)$){#4};
\node[right] at ($(B)!0.5!(C)$){#3};
\node[above left] at ($(A)!0.5!(C)$){#5};
\node at (0.7,0.22){$\theta$};
\end{tikzpicture}\end{center}}

% General (oblique) triangle with vertex labels and opposite-side labels:
%   \gtri{Alabel}{Blabel}{Clabel}{aSide}{bSide}{cSide}
\newcommand{\gtri}[6]{\begin{center}\begin{tikzpicture}[scale=0.85]
\coordinate (A) at (0,0);\coordinate (B) at (5,0);\coordinate (C) at (1.6,3);
\draw[thick,exblue] (A)--(B)--(C)--cycle;
\node[below left] at (A){#1};\node[below right] at (B){#2};\node[above] at (C){#3};
\node[below] at ($(A)!0.5!(B)$){#6};
\node[right] at ($(B)!0.5!(C)$){#4};
\node[left] at ($(A)!0.5!(C)$){#5};
\end{tikzpicture}\end{center}}

\fancyhead[R]{\footnotesize MCV4U \textperiodcentered{} 6.4}
\begin{document}

\begin{center}
{\LARGE\bfseries\color{titleblue}Worksheet 6.4 --- The Dot Product}\\[2pt]
{\small Grade 12 Calculus and Vectors (MCV4U) \textperiodcentered{} Unit 6: Geometry and Algebra of Vectors}
\end{center}

\begin{ideabox}
$\vec a\cdot\vec b=a_1b_1+a_2b_2(+a_3b_3)=|\vec a||\vec b|\cos\theta$. Gives the angle; zero means perpendicular; work $=\vec F\cdot\vec d$.
\begin{itemize}[leftmargin=*,itemsep=2pt]
  \item Components: $\vec a\cdot\vec b=a_1b_1+a_2b_2+a_3b_3$.
  \item $\cos\theta=\dfrac{\vec a\cdot\vec b}{|\vec a||\vec b|}$.
  \item Perpendicular $\iff\vec a\cdot\vec b=0$; work $=\vec F\cdot\vec d$.
\end{itemize}
\end{ideabox}

\begin{learnbox}
\lh{Multiplying to a scalar}The \textbf{dot product} $\vec u\cdot\vec v=u_1v_1+u_2v_2+\cdots$ produces a \emph{scalar}, not a vector.
\lh{The angle connection}It also equals $|\vec u|\,|\vec v|\cos\theta$, so it measures how aligned two vectors are and gives the angle between them.
\lh{Perpendicular test}Two nonzero vectors are perpendicular exactly when their dot product is $0$.
\end{learnbox}

\sech{Worked Examples}

\begin{example}{ex0}{Example 1 --- Dot product (2-D)}
Find $(2,3)\cdot(4,-1)$.\soln
\textbf{Step 1 (multiply matching components):} $(2)(4)$ and $(3)(-1)$.\\[3pt]
\textbf{Step 2 (add the products):} $8+(-3)$.\\[3pt]
\textbf{Conclusion:} $=5$.
\end{example}

\begin{example}{ex1}{Example 2 --- Dot product (3-D)}
Find $(1,2,3)\cdot(2,0,-1)$.\soln
\textbf{Step 1 (multiply matching components):} $(1)(2)$, $(2)(0)$, $(3)(-1)$.\\[3pt]
\textbf{Step 2 (add the products):} $2+0+(-3)$.\\[3pt]
\textbf{Conclusion:} $=-1$.
\end{example}

\begin{example}{ex2}{Example 3 --- Angle}
Angle between $(1,0)$ and $(1,1)$?\soln
\textbf{Step 1 (find the dot product):} $(1)(1)+(0)(1)=1$.\\[3pt]
\textbf{Step 2 (find both magnitudes):} $|(1,0)|=1$; $|(1,1)|=\sqrt2$.\\[3pt]
\textbf{Step 3 (apply $\cos\theta=\frac{\vec a\cdot\vec b}{|\vec a||\vec b|}$):} $\cos\theta=\dfrac{1}{1\cdot\sqrt2}=\dfrac1{\sqrt2}$.\\[3pt]
\textbf{Conclusion:} $\theta=45^\circ$.
\end{example}

\begin{example}{ex3}{Example 4 --- Perpendicular}
Are $(2,3)$ and $(3,-2)$ perpendicular?\soln
\textbf{Step 1 (compute the dot product):} $(2)(3)+(3)(-2)=6-6$.\\[3pt]
\textbf{Step 2 (recall the perpendicularity test, dot product $=0$):} the result is $0$.\\[3pt]
\textbf{Conclusion:} yes, perpendicular.
\end{example}

\begin{example}{ex4}{Example 5 --- Work}
$\vec F=(3,4)$ N, $\vec d=(2,0)$ m. Find the work.\soln
\textbf{Step 1 (recall the work formula):} $W=\vec F\cdot\vec d$.\\[3pt]
\textbf{Step 2 (compute the dot product):} $(3)(2)+(4)(0)=6+0$.\\[3pt]
\textbf{Conclusion:} $W=6$ J.
\end{example}

\begin{example}{ex5}{Example 6 --- Dot product}
Find $(1,2)\cdot(3,4)$.\soln
\textbf{Step 1 (multiply matching components and add):} $(1)(3)+(2)(4)=3+8$.\\[3pt]
\textbf{Conclusion:} $=11$.
\end{example}

\begin{example}{ex6}{Example 7 --- Dot product (3-D)}
Find $(2,-1,3)\cdot(1,4,2)$.\soln
\textbf{Step 1 (multiply matching components and add):} $(2)(1)+(-1)(4)+(3)(2)=2-4+6$.\\[3pt]
\textbf{Conclusion:} $=4$.
\end{example}

\begin{example}{ex7}{Example 8 --- Angle}
Angle between $(1,0)$ and $(0,1)$?\soln
\textbf{Step 1 (compute the dot product):} $(1)(0)+(0)(1)=0$.\\[3pt]
\textbf{Step 2 (interpret a zero dot product):} the vectors are perpendicular.\\[3pt]
\textbf{Conclusion:} $90^\circ$.
\end{example}

\begin{example}{ex8}{Example 9 --- Work}
$\vec F=(5,0)$, $\vec d=(3,4)$. Find the work.\soln
\textbf{Step 1 (apply $W=\vec F\cdot\vec d$):} $(5)(3)+(0)(4)=15+0$.\\[3pt]
\textbf{Conclusion:} $W=15$ J.
\end{example}

\clearpage
\sech{Your Turn --- Show Your Work}
For each question, show all steps clearly in the space provided.

\begin{qbox}{Question 1}
Find $(1,2)\cdot(3,4)$.\workspace{2.4cm}
\end{qbox}

\begin{qbox}{Question 2}
Find $(2,-1,3)\cdot(1,4,2)$.\workspace{2.4cm}
\end{qbox}

\begin{qbox}{Question 3}
Are $(1,2)$ and $(4,-2)$ perpendicular?\workspace{2.4cm}
\end{qbox}

\begin{qbox}{Question 4}
Angle between $(1,0)$ and $(0,1)$?\workspace{2.4cm}
\end{qbox}

\begin{qbox}{Question 5}
Work of $\vec F=(5,0)$ along $\vec d=(3,4)$?\workspace{2.4cm}
\end{qbox}

\begin{qbox}{Question 6}
Find $(3,1)\cdot(2,5)$.\workspace{2.4cm}
\end{qbox}

\begin{qbox}{Question 7}
Find $(1,1,1)\cdot(2,3,4)$.\workspace{2.4cm}
\end{qbox}

\begin{qbox}{Question 8}
Are $(3,4)$ and $(4,-3)$ perpendicular?\workspace{2.4cm}
\end{qbox}

\begin{qbox}{Question 9}
Angle between $(2,0)$ and $(2,2)$?\workspace{2.4cm}
\end{qbox}

\begin{qbox}{Question 10}
Work of $\vec F=(2,3)$ along $\vec d=(4,0)$?\workspace{2.4cm}
\end{qbox}

\begin{qbox}{Question 11}
Find $(0,5)\cdot(5,0)$.\workspace{2.4cm}
\end{qbox}

\begin{qbox}{Question 12}
Find $|\vec a|$ from $\vec a\cdot\vec a=25$.\workspace{2.4cm}
\end{qbox}

\begin{qbox}{Question 13 --- Challenge}
Find the angle between $(3,4)$ and $(4,3)$ to the nearest degree.\workspace{3cm}
\end{qbox}

\clearpage
\sech{Question Recap \& Answer Key}

\begin{recapbox}
\begin{multicols}{2}
\small
\begin{enumerate}[leftmargin=*,itemsep=3pt]
  \item Find $(1,2)\cdot(3,4)$.
  \item Find $(2,-1,3)\cdot(1,4,2)$.
  \item Are $(1,2)$ and $(4,-2)$ perpendicular?
  \item Angle between $(1,0)$ and $(0,1)$?
  \item Work of $\vec F=(5,0)$ along $\vec d=(3,4)$?
  \item Find $(3,1)\cdot(2,5)$.
  \item Find $(1,1,1)\cdot(2,3,4)$.
  \item Are $(3,4)$ and $(4,-3)$ perpendicular?
  \item Angle between $(2,0)$ and $(2,2)$?
  \item Work of $\vec F=(2,3)$ along $\vec d=(4,0)$?
  \item Find $(0,5)\cdot(5,0)$.
  \item Find $|\vec a|$ from $\vec a\cdot\vec a=25$.
  \item Find the angle between $(3,4)$ and $(4,3)$ to the nearest degree.
\end{enumerate}
\end{multicols}
\end{recapbox}

\begin{keybox}
\begin{multicols}{2}
\small
\begin{enumerate}[leftmargin=*,itemsep=3pt]
  \item $11$
  \item $4$
  \item $0$: yes
  \item $90^\circ$
  \item $15$ J
  \item $11$
  \item $9$
  \item $0$: yes
  \item $45^\circ$
  \item $8$ J
  \item $0$
  \item $5$
  \item $\cos\theta=\tfrac{24}{25}\Rightarrow\theta\approx16^\circ$
\end{enumerate}
\end{multicols}
\end{keybox}

\end{document}
